\chapter*{Prólogo}
\addcontentsline{toc}{chapter}{Prólogo}

Este volumen cubre la asignatura Mecánica de Fluidos del Grado en Ingeniería Aeroespacial de la Universidad de León, código 0710311, según la guía docente firmada para 2026/2027 \parencite{uleGuia2026,ulePlan2026}. Está escrito para estudiar la materia, no para sustituir un formulario.

Cada resultado importante lleva sus hipótesis. Una ecuación sin ese marco no es una ley del curso: es una frase a medias. Las derivaciones se hacen en el texto. Cuando un paso se deja al lector, el paso no cambia el resultado.

La guía ordena los casos antes que parte del formalismo. Este libro invierte ese tramo a propósito: primero el lenguaje del continuo, la cinemática y los balances, y después Couette, Poiseuille, la resistencia y las aplicaciones. El mapa entre los seis temas oficiales y los capítulos está en el capítulo~\ref{ch:continuo}.

\begin{warningbox}
La guía vigente no incluye cohetes ni una atmósfera estándar. Esos temas aparecen en guías anteriores y en material de estudiantes. No forman parte de esta edición. Los enunciados de examen ajenos tampoco se reproducen: los problemas de aquí son nuevos y buscan las mismas habilidades.
\end{warningbox}

Las lecturas que la guía recomienda se citan como tales. No se transcriben. Los resultados matemáticos de este volumen se sostienen en las derivaciones que el lector puede repetir.
